%---------------------------------------
% References
%---------------------------------------
\clearpage
\phantomsection
\addcontentsline{toc}{chapter}{References}
\bibliographystyle{plain}
\bibliography{references}

%---------------------------------------
% Appendices
%---------------------------------------
\appendix

\chapter{Hardware Schematics and Pinouts}
\label{app:hardware}

This appendix details the physical wiring between the microcontroller and the measurement instrumentation described in Section~\ref{sec:hil-architecture}.

\begin{figure}[H]
\centering
\begin{tikzpicture}[
    node distance=1.2cm and 3.0cm,
    box/.style={draw, rounded corners, align=center, minimum height=1.2cm, minimum width=3.2cm, font=\small},
    inst/.style={draw, rounded corners, align=center, minimum height=1.2cm, minimum width=3.2cm, font=\small, fill=gray!12},
    arr/.style={-{Latex[length=2mm]}, thick}
]
\node[box] (mcu) {Microcontroller\\(ESP32 / Teensy)};
\node[inst, right=of mcu] (la) {Logic analyzer /\\oscilloscope};
\node[box, below=1.6cm of mcu] (host) {Host PC (UART/USB)};

\draw[arr] (mcu.north east) -- node[midway, above, font=\scriptsize]{GPIO\,A -- \texttt{t0\_start}} (la.north west);
\draw[arr] (mcu.south east) -- node[midway, below, font=\scriptsize]{GPIO\,B -- \texttt{spike\_out}} (la.south west);
\draw[arr] (host) -- node[left, font=\scriptsize, align=center]{serial / SPI\\(input spike train)} (mcu);
\node[font=\scriptsize, below=0.1cm of host] {GND -- common reference shared with all instruments};
\end{tikzpicture}
\caption[Hardware wiring schematic]{Wiring schematic between the host PC, the microcontroller, and the logic analyzer/oscilloscope. \texttt{t0\_start} marks the beginning of inference; \texttt{spike\_out} marks the first output spike (physical TTFS event). A common ground reference is shared by all instruments.}
\label{fig:hw-schematic}
\end{figure}

\begin{table}[H]
\centering
\caption[Pinout summary]{Summary of instrumented microcontroller pins.}
\label{tab:pinout}
\small
\begin{tabular}{@{}llp{7.2cm}@{}}
\toprule
\textbf{Signal} & \textbf{Direction} & \textbf{Function} \\
\midrule
\texttt{t0\_start} & Output & Rising edge at the start of inference (timestep $t=1$) \\
\texttt{spike\_out} & Output & Rising edge on first spike of the winning output neuron \\
UART/SPI RX & Input  & Reception of the precomputed spike-encoded input tensor \\
GND         & --     & Common reference shared with the logic analyzer/oscilloscope \\
\bottomrule
\end{tabular}
\end{table}

\chapter{Embedded Inference Code}
\label{app:code}

This appendix reports the most critical fragments of the embedded firmware described in Section~\ref{sec:porting}: the quantized LIF membrane update and the GPIO-toggling inference loop.

\begin{lstlisting}[caption={Quantized LIF membrane update (fixed-point arithmetic).}, label={lst:qlif}]
// Quantized Leaky Integrate-and-Fire update.
// weights, mem, and beta_q are pre-quantized integers exported
// from the QAT model; shift implements the fixed-point
// scale factor of the membrane decay.
int32_t update_lif_neuron(int32_t mem, int32_t input_current,
                           int32_t beta_q, int32_t shift,
                           int32_t v_th_q, uint8_t *spike_out) {
    // Leak: mem = beta_q * mem (fixed-point, right-shifted)
    mem = (int32_t)(((int64_t)beta_q * mem) >> shift);

    // Integrate weighted input current for this timestep
    mem += input_current;

    // Fire and soft-reset (subtract threshold, retain residual)
    if (mem >= v_th_q) {
        *spike_out = 1;
        mem -= v_th_q;
    } else {
        *spike_out = 0;
    }
    return mem;
}
\end{lstlisting}

\begin{lstlisting}[caption={GPIO-instrumented inference loop for the Hardware-in-the-Loop protocol (Section~\ref{sec:measurement-protocol}).}, label={lst:gpio-loop}]
#define PIN_T0_START   GPIO_NUM_A
#define PIN_SPIKE_OUT  GPIO_NUM_B

void run_hil_inference(const uint8_t spike_train[T][N_IN]) {
    reset_network_state();               // clear all membrane potentials

    gpio_set_level(PIN_T0_START, 1);      // t = 0 reference edge
    gpio_set_level(PIN_T0_START, 0);

    for (int t = 0; t < T; t++) {
        forward_step(spike_train[t]);     // quantized conv/fc + LIF layers

        uint8_t winner_spiked = 0;
        int winner_class = argmax_readout_spike(&winner_spiked);

        if (winner_spiked && !decision_latched) {
            gpio_set_level(PIN_SPIKE_OUT, 1);   // physical TTFS event
            gpio_set_level(PIN_SPIKE_OUT, 0);
            decision_latched = 1;               // record t* implicitly
        }
    }
}
\end{lstlisting}
